\documentclass[12pt]{article}
\usepackage[utf8]{inputenc}
\usepackage[T1]{fontenc}
\usepackage[spanish]{babel}
\usepackage[margin=2.5cm]{geometry}
\usepackage{amsmath}
\usepackage{enumitem}

\pagestyle{empty}
\setlength{\parindent}{0pt}
\setlength{\parskip}{6pt}

\begin{document}

{\large\textbf{Econometría II · Hoja de ejercicios del domingo 1-11}}

\textbf{Ejercicio 1. Logit sin mirar (30 min)}

{\small Ejercicio inventado para practicar: no viene de una prueba de Pfeng. La ecuación es la de la clase 3 (Spector y Mazzeo); el perfil y las preguntas son nuevos, y la definición de $Y$ («mejora su calificación») falta confirmarla con los apuntes.}

Resuelve en papel, con la calculadora de la prueba y con cronómetro. Entrega cada resultado con su frase de interpretación completa: perfil, unidad, «aproximadamente» y «manteniendo lo demás constante» cuando corresponda.

\[
\hat{z} = -13{,}0213 + 2{,}8261\,\text{GPA} + 0{,}0952\,\text{TUCE} + 2{,}3787\,\text{PSI},
\qquad \text{con } \text{GPA}=3{,}0 \text{ y } \text{TUCE}=25.
\]

\begin{enumerate}[label=\alph*), itemsep=4pt]
  \item $\hat{P}$ con $\text{PSI}=0$ y con $\text{PSI}=1$.
  \item Efecto de PSI e interpretación.
  \item Efecto marginal de GPA con $\text{PSI}=0$, interpretación y por qué no es $2{,}8261$. Con ese efecto marginal, ¿cuál es el efecto de un aumento de $0{,}1$ en GPA?
  \item Cambio en los odds por PSI.
  \item Clase estimada con corte $0{,}5$.
  \item Dos razones para preferir logit al MPL.
\end{enumerate}

\end{document}
